% =========================================================================
% anexos/01_anexos.tex
% Sección de Anexos del Informe
% =========================================================================

\subsection{Especificación Técnica y Diccionario de Datos}
\label{anexo:diccionario_datos}

En este anexo se pueden incluir tablas técnicas extensas, diagramas de base de datos o especificaciones de interfaces que por su extensión no deben interrumpir la lectura del cuerpo principal.

\begin{table}[htbp]
  \centering
  \caption{Estructura representativa de la entidad principal de datos.}
  \label{tab:anexo_entidad_datos}
  \begin{tabularx}{\textwidth}{llLX}
    \toprule
    \textbf{Campo} & \textbf{Tipo} & \textbf{Restricción} & \textbf{Descripción} \\
    \midrule
    \texttt{id} & \texttt{UUID} & Primary Key & Identificador único inmutable del registro. \\
    \texttt{codigo} & \texttt{VARCHAR(20)} & Unique, Not Null & Código administrativo institucional. \\
    \texttt{estado} & \texttt{VARCHAR(15)} & Not Null & Estado operativo (Activo, Pendiente, Inactivo). \\
    \texttt{fecha\_creacion} & \texttt{TIMESTAMP} & Default NOW() & Registro cronológico de auditoría. \\
    \bottomrule
  \end{tabularx}
  \fuente{Elaboración propia.}
\end{table}

\subsection{Guía de Despliegue y Scripts de Automatización}
\label{anexo:scripts_despliegue}

En este anexo se documentan los pasos o fragmentos de código utilizados para el aprovisionamiento de entornos:

\begin{lstlisting}[language=bash, caption={Script de automatización de pruebas y compilación del proyecto.}, label={lst:script_automatizacion}]
#!/usr/bin/env bash
# Script para compilación y verificación automatizada
set -e

echo "Iniciando pruebas unitarias..."
pytest --cov=src tests/

echo "Compilando contenedor de la aplicación..."
docker build -t empresa/sistema-gestion:latest .

echo "Despliegue completado con éxito."
\end{lstlisting}
